\documentclass[10pt,a4paper]{amsart}
\usepackage[margin=19mm]{geometry}
\usepackage{amsmath,amssymb,mathtools}
\usepackage[scr=rsfs,cal=euler]{mathalfa}
\usepackage{xfrac}
\usepackage[unicode,hidelinks,bookmarks=false]{hyperref}
\newcommand{\ie}{i.e.\ }
\newcommand{\cf}{cf.\ }
\newcommand{\resp}{respectively\ }
\newcommand{\Subsection}{\subsection}
\providecommand{\nobreakdash}{\nobreak\hbox{-}}
\setlength{\emergencystretch}{2em}
\multlinegap0pt
\usepackage{bm}
\usepackage{bbm}
\usepackage{dsfont}
\usepackage{xspace}
\usepackage{cancel}
\usepackage{mathtools}
\usepackage{amsthm}
\makeatletter
\@ifundefined{theorem}{\newtheorem{theorem}{Theorem}}{}
\@ifundefined{lemma}{\newtheorem{lemma}[theorem]{Lemma}}{}
\makeatother
\title[Double signs of Hamiltonian circles in doubly signed complet]{Double signs of Hamiltonian circles in doubly signed complete graphs}
\author{Xiyong Yan}
\date{Source: arXiv:2509.16646v1 (CC BY 4.0)}
\begin{document}
\maketitle

\noindent\emph{Source and licence.} Double signs of Hamiltonian circles in doubly signed complete graphs. Version arXiv:2509.16646v1 (CC BY 4.0), distributed under the Creative Commons Attribution 4.0 International licence, as recorded in the arXiv licence metadata for this exact version and shown on the abstract page https://arxiv.org/abs/2509.16646v1. Full rights notice: licenses/arxiv-2509.16646-CC-BY-4.0.txt.\par\smallskip

\noindent\textit{Editorial scope.} Counted unit: the Lemma whose source label is \texttt{lemma4} in the article (source0.tex lines 967--972, source label \texttt{lemma4}), delivered together with the article's own complete original proof of it (source0.tex lines 974--1024, one \texttt{proof} environment of 51 lines). No mathematics is written, repaired or re-derived: statement and proof are the article's own text line for line. Required context retained uncounted: none. Every definition, display and label of those lines is retained unchanged. Omitted from this package and not redistributed: 1--966, 973--973, 1025--1566. Nothing inside a retained excerpt is omitted or altered: every retained body is delivered exactly as the article's lines are written, so no omission receipt of any kind accompanies this package. Citations: the retained excerpts carry no citation command at all, so no bibliography entry of the article is transcribed and no reference is redistributed. Reference adaptation: none was needed and none was made; every reference inside the delivered unit resolves inside it, so no unresolved reference or citation is produced. Printed slips kept literal: no printed word, symbol, hypothesis or number of the counted statement or of its proof was altered. Layout and class: the article's own class is \texttt{amsart} with packages including graphicx, tikz, xcolor, amsmath, caption, geometry, float, mathrsfs, fancyhdr, amssymb; the wrapper is the lane's pinned \texttt{amsart} editorial wrapper at 10pt on A4, which restates the theorem-like environments the retained text needs and adds the article's own macro definitions that the retained text needs (none), with extra wrapper packages (bm, bbm, dsfont, xspace, cancel, mathtools). The delivered numbering is the unit's own. Assets: no retained span contains a figure environment, a graphics-inclusion command, a PDF-page inclusion, a table or a TikZ picture; no image file is copied into this package, the package declares no asset dependency and no image file is redistributed with it. Licence: version arXiv:2509.16646v1 of this article is distributed under the Creative Commons Attribution 4.0 International licence, as recorded in the arXiv licence metadata for this exact version and shown on the abstract page https://arxiv.org/abs/2509.16646v1. A full rights notice is delivered at licenses/arxiv-2509.16646-CC-BY-4.0.txt. Characters: every retained excerpt is pure ASCII, so the delivered engine, XeLaTeX with its default Latin Modern text font, reports no missing character.
\begin{lemma}\label{lemma4}
For each \(i\in\{1,2,3,4\}\), the multiset of double signs of the six Hamiltonian paths
\(\mathscr{P}_i\) starting at \(v_i\) is either
\(\{p,p,q,q,s,s\}\) or \(\{p,q,s,t,t,t\}\) for some distinct
\(p,q,s,t\in\mathbb F_2^2\).
\end{lemma}

\begin{proof}
By symmetry it suffices to take \(i=1\).
Assume (WLOG for \(\Sigma_4^*\)) that
\(\sigma(T_{123})=a\), \(\sigma(T_{124})=b\), \(\sigma(T_{134})=c\), and
\(\sigma(T_{234})=e\) (the identity). 

Let \(x=\sigma(e_{12})\), \(y=\sigma(e_{13})\), \(z=\sigma(e_{14})\).
The six Hamiltonian paths from \(v_1\) arise by deleting one of the three
edges \(e_{12},e_{13},e_{14}\) from one of the Hamiltonian cycles
\(C_1=v_1v_2v_3v_4\), \(C_2=v_1v_2v_4v_3\), \(C_3=v_1v_3v_2v_4\).
Since \(\sigma(C_1)=b\), \(\sigma(C_2)=a\), \(\sigma(C_3)=c\) and $\sigma(C_i\backslash e_{jk})=\sigma(C_i)+\sigma(e_{jk})$,
the six paths   \(C_{2}\setminus e_{12}\), \(C_{1}\setminus e_{12}\),
\(C_{2}\setminus e_{13}\),  \(C_{3}\setminus e_{13}\), $C_1\backslash e_{14},$ and $ C_3\backslash e_{14}$ 
have signs 
\[
\underbrace{a+x,\ b+x}_{\text{omit }e_{12}},\qquad
\underbrace{a+y,\ c+y}_{\text{omit }e_{13}},\qquad
\underbrace{b+z,\ c+z}_{\text{omit }e_{14}},
\] respectively.

(1) Within each brace, the two values are distinct.
Indeed,
\((a+x)+(b+x)=a+b=c\neq e\), and \((b+z)+(c+z)=b+c=a\neq e\), hence
\(a+x\ne b+x\) and \(b+z\ne c+z\). Similarly \((a+y)+(c+y)=a+c=b\ne e\), so
\(a+y\ne c+y\).

(2) Between each pair of braces, four double signs cannot be all distinct. 
For example, between \(\{a+x,b+x\}\) and \(\{b+z,c+z\}\):  Suppose all of them are different. $a+x\neq b+z$ or $c+z$ implies $x\neq c+z$ or $b+z$. $b+x\neq b+z$ or $c+z$ implies $x\neq e+z$ or $a+z$. It would force 
$x$ to avoid all four elements of $\mathbb{F}_2^2$, a contradiction. Thus, $a+x$, $b+x$, $b+z$ and $c+z$ cannot be all different.


(3) Between each pair of braces, exactly one element from one brace equals an element from the other.
 Suppose there are two pairs of them have the same double signs. Assume 
     $a+x=b+z$, $b+x=c+z$. This implies $x=c+z$ and $x=a+z$, impossible. Similary, $a+x=c+z$, $b+x=b+z$ cannot happen. Hence, there are exactly two of $a+x$, $b+x$, $b+z$ and $c+z$ are the same. 

\vspace{0.5cm}

Finally, we can analyze the double signs of the six paths starting at $v_1$.

Case 1. Three of $a+y,c+y$, $a+x$, $b+x$, $b+z$ and $c+z$ are the same. Then, remaining three of them are all different. In this case, we have all four double signs $a,b,c,e$ with exactly one of them appearing three times. 

Case 2. No three elements of $a+y,c+y$, $a+x$, $b+x$, $b+z$ and $c+z$  share the same double sign. Consequently, there are
precisely three pairs of equal values, and the value of each pair is
different from the others.

\vspace{0.5cm}

Therefore, \(\sigma(\mathscr P_1)\) is either \(\{p,p,q,q,s,s\}\) or
\(\{p,q,s,t,t,t\}\) with \(p,q,s,t\) distinct. By symmetry the same holds for
\(i=2,3,4\).
\end{proof}

\end{document}
